% ============================================================
% 00-abstract.tex
% ============================================================
\begin{abstract}
Student counseling in higher education draws on heterogeneous institutional knowledge, including relational curricula, cohort-specific tuition schedules, statutory exemption criteria, scholarship policies, and narrative administrative regulations. Monolithic RAG systems and loosely coordinated agent architectures can be unreliable in this setting because different forms of knowledge require different retrieval and reasoning methods. We present CTU-Chat, a centralized supervisor-routed, domain-specialized multi-agent system for Vietnamese university counseling. The supervisor directs each query to one of four specialists, while each specialist is assigned an appropriate knowledge source, retrieval path, and set of tools. Relational academic information is handled through a knowledge graph; structured tuition and scholarship information is supported by dedicated lookup and calculation functions; and narrative regulations are retrieved through hybrid document search. This division allows the system to use each source according to its structure while limiting unnecessary tool access. Evaluation shows that CTU-Chat achieves 97.0\% supervisor routing accuracy, while its proposed retrieval configuration reaches 0.9600 Hit@1 on the held-out test set. These findings suggest that combining centralized coordination with domain-specific knowledge allocation can provide a practical and reliable foundation for university counseling systems.

\keywords{Multi-Agent Systems \and Domain-Specialized Agents \and Centralized Supervisor Routing \and Heterogeneous Knowledge Allocation \and Retrieval-Augmented Generation \and University Counseling}
\end{abstract}
